\documentclass[../main]{subfiles}
\begin{document}

\section{Análisis de Criticidad}
\img{images/07/logo_criticidad}{Análisis de Criticidad}
\info{Definición de análisis de criticidad}{El análisis de criticidad es una metodología que permite
establecer la jerarquía o prioridades de procesos, sistemas y equipos, creando una estructura que facilita la toma de
decisiones acertadas y efectivas, direccionando el esfuerzo
y los recursos en áreas donde sea más importante y/o
necesario mejorar la fiabilidad operacional, basado en la
realidad actual.
La mejora de la fiabilidad operacional de cualquier
instalación o de sus sistemas y componentes, está asociado
con cuatro aspectos fundamentales: fiabilidad del proceso,
fiabilidad humana, fiabilidad de los equipos y
mantenimiento de los equipos.}
\subsection{Proceso de análisis de criticidad.}
Los pasos a realizar un análisis de criticidad son:
\begin{enumerate}
    \item Identificación de equipos a estudiar y realizar una lista de equipos.
\item Definición de alcance y objetivos del análisis,
jerarquía, nivel y profundidad al que se quiere
aplicar el estudio.
\item Selección del grupo de trabajo. Encargado de
definir el impacto de producción, económico, de
personal, ambiental, costes de mantenimiento y
seguridad.
\item Interpretación de datos y aplicación de la tabla de criticidad. 
\end{enumerate}

\imgfull{images/07/tabla_criticidad}{Tabla de análisis de criticidad}

\subsection{Niveles de importancia}
No todos los equipos tienen la misma importancia en una planta industrial. Es un hecho que unos equipos son más importantes que otros. Como los recursos de una empresa para mantenerla son limitados, se debe destinar a la mayor parte de los recursos a los equipos más importantes, la pequeña parte del reparto a los equipos que menos pueden influir en los resultados de una empresa.

\begin{enumerate}
    \item \emph{Equipos críticos:} Aquellos de los que dependen los resultados de la empresa.
\item \emph{Equipos importantes:} Afectan los objetivos de la empresa, pero las consecuencias son asumibles.
\item \emph{Equipos prescindibles:} Son aquellos con una incidencia en los resultados mínima. Como mucho, supongamos una pequeña incomodidad, un pequeño cambio de trascendencia, o un pequeño costo adicional.
\end{enumerate}
\info{Equipos altamente críticos}{Opcionalmente, algunas empresas prefieren incluir una categoría más: los equipos altamente críticos. Se pretende con la introducción de esta nueva categoría.
}
\subsection{Consecuencias que afectan en el análisis}
 Se debe considerar la influencia que una anomalía tiene cuatro aspectos: producción, calidad, mantenimiento y seguridad.
\begin{enumerate}
    \item  Producción:  Cuando se valora la influencia que tiene un equipo en la producción, se debe preguntar cómo hacerlo esta vez una posible falla. Dependiendo de que supongamos una parada total de la instalación, una parada de una zona de producción preferente, paralizar los equipos productivos pero los menores producciones no se puede influir en la producción, se clasifica como el equipo como A, B o C.

    \item  Calidad: El equipo puede tener una influencia decisiva en la calidad del producto o servicio final, una influencia relativa que no puede ser una tarea problemática o una influencia nula. 

    \item  Mantenimiento: El equipo puede ser muy problemático, con averías caras o frecuentes; o bien un equipo con un costo medio en mantenimiento; o por último, un equipo con un costo muy bajo. 

    \item Seguridad y medio ambiente:  Una falla del equipo puede significar un accidente muy grave, bien por el tiempo o por las personas, y también tener una probabilidad de falla; es posible que una falla del equipo pueda ocasionar un accidente, pero la probabilidad de que esto ocurra puede ser baja; o por último, puede ser un equipo que no tenga ninguna influencia en la seguridad. 
\end{enumerate}

 


\subsection{Diagrama de Pareto}
\img{images/07/relleno1}{Análisis de datos.}
Según el principio de Pareto adaptado a la gestión de
calidad por Joseph M. Juran (1904-2008).
Aproximadamente el 80\% de las variaciones que se dan en
un proceso, están originadas por el 20\% de las causas de variación presentes en dicho proceso.
Consecuentemente si resolvemos el 20\% de las causas de
variación, estaremos eliminando el 80\% de la variación
existente en el proceso. El diagrama de Pareto consiste en un gráfico de barras que muestra ordenadas de mayor a
menor las frecuencias de determinados hechos


\newpage
\actividad{Leer y responder}
\begin{enumerate}
    \item ¿En qué consiste el análisis de criticidad?
    \item ¿Qué es la fiabilidad?
    \item ¿Qué se entiende por Calidad?
    \item ¿Qué es una externalidad y como afectan las externalidades al medio ambiente?
    \item ¿Cómo identificaría de forma rápida a los equipos prescindibles en función de que factores?
\end{enumerate}

\actividad{Investigar}
\begin{enumerate}
    \item ¿Cómo se diferencian los equipos que tiene una gran influencia en los resultados de los que no lo tienen? 
    \item ¿Quién fué Pareto, y qué aplicaciones tiene el diagrama de Pareto?
    
\end{enumerate}
\end{document}
